\begin{pastexam}[label={px:ntu-110-ie-ds-21}]{臺大工工所 110 學年度 第 21 題}
（複選；每題 5 分，每答錯一個選項倒扣 1 分至該題 0 分止）

Consider an AA tree that results from successively inserting keys 96, 49, 79, 14, 41, 87, 75, 90, 61 into
an initially empty tree. Which of the following statements are true?
\begin{enumerate}[label=(\Alph*)]
\item The root is key 61.
\item There is less than 3 horizontal links.
\item Key 41 and key 49 are in the same level.
\item Key 87 and key 96 are siblings.
\item Key 14 is the parent of key 41.
\end{enumerate}
\end{pastexam}

\begin{pastsoln}
答案是 (B)、(D)。

AA 樹資結書沒有專節，先說明規則（Andersson 1993，Weiss 的教科書用的版本）。 每個節點有一個\textbf{層級}（level），葉節點是 $1$：
\begin{enumerate}
\item 左子節點的層級比父節點小 $1$；
\item 右子節點的層級等於父節點或小 $1$（相等時這條邊叫\textbf{水平連結}，相當於紅黑樹的紅色右子節點）；
\item 右孫節點的層級一定比祖父節點小（不能連續兩條水平連結）；
\item 層級大於 $1$ 的節點有兩個子節點。
\end{enumerate}
插入時照二元搜尋樹放成層級 $1$ 的葉節點，再沿著路徑往回，每個節點先做 \textbf{skew}、再做 \textbf{split}：
\begin{itemize}
\item skew：左子節點和自己同層（出現往左的水平連結）時右旋，讓水平連結改成往右。
\item split：右子節點和右孫節點都和自己同層（連續兩條水平連結）時左旋，讓中間那個節點升一層。
\end{itemize}
下面 $k/\ell$ 表示鍵 $k$、層級 $\ell$。
\begin{enumerate}
\item $96$：單獨一個節點 $96/1$。
\item $49$：放在 $96$ 左邊；在 $96$ 做 skew，變成 $49/1$，右子節點 $96/1$（水平）。
\item $79$：路徑 $49 \to 96$，放在 $96$ 左邊；在 $96$ skew，變成 $79 \to 96$ 往右水平；回到 $49$，$49 \to 79 \to 96$ 連續兩條水平連結，
split：$79$ 升為 $79/2$，左 $49$、右 $96$。
\item $14$：路徑 $79 \to 49$，放在 $49$ 左邊；在 $49$ skew，變成 $14 \to 49$（水平）。
\item $41$：路徑 $79 \to 14 \to 49$，放在 $49$ 左邊；在 $49$ skew 成 $41 \to 49$；回到 $14$，$14 \to 41 \to 49$ 連續水平，split，$41$ 升為
$41/2$（左 $14$、右 $49$）；回到 $79$，左子節點 $41$ 和 $79$ 同為 $2$ 層，skew：$41$ 成為根，右子節點 $79/2$（水平），$49$ 改成
$79$ 的左子節點。 樹：$41/2$，左 $14/1$，右 $79/2$（左 $49/1$、右 $96/1$）。
\item $87$：路徑 $41 \to 79 \to 96$，放在 $96$ 左邊；skew 成 $87 \to 96$（水平）。 $79$ 的右子節點 $87$ 在第 $1$ 層，不必 split。
\item $75$：路徑 $41 \to 79 \to 49$，放在 $49$ 右邊（水平），不必調整。
\item $90$：路徑 $41 \to 79 \to 87 \to 96$，放在 $96$ 左邊；skew 成 $90 \to 96$；回到 $87$，$87 \to 90 \to 96$ 連續水平，split，$90$ 升為
$90/2$（左 $87$、右 $96$）；回到 $79$，右子節點 $90/2$ 和它同層，是水平連結；回到 $41$，$41 \to 79 \to 90$ 三個都在第 $2$ 層，
split，$79$ 升為 $79/3$ 成為根，左 $41/2$、右 $90/2$。
\item $61$：路徑 $79 \to 41 \to 49 \to 75$，放在 $75$ 左邊；skew 成 $61 \to 75$；回到 $49$，$49 \to 61 \to 75$ 連續水平，split，$61$ 升為
$61/2$（左 $49$、右 $75$）；回到 $41$，右子節點 $61/2$ 和它同層，是一條水平連結，$61$ 的右子節點 $75$ 在第 $1$ 層，不必 split。
\end{enumerate}
最後的樹（虛線是水平連結，同一層的節點畫在同一高度）：
\begin{center}
\begin{tikzpicture}[x=1cm,y=1.1cm,v/.style={circle,draw,minimum size=6.5mm,inner sep=0pt,font=\small}]
\node[v] (79) at (3.3,0) {79};
\node[v] (41) at (0.9,-1) {41}; \node[v] (61) at (2.4,-1) {61}; \node[v] (90) at (5.4,-1) {90};
\node[v] (14) at (0,-2) {14}; \node[v] (49) at (1.7,-2) {49}; \node[v] (75) at (3.1,-2) {75};
\node[v] (87) at (4.7,-2) {87}; \node[v] (96) at (6.1,-2) {96};
\draw (79)--(41); \draw (79)--(90); \draw[dashed] (41)--(61);
\draw (41)--(14); \draw (61)--(49); \draw (61)--(75); \draw (90)--(87); \draw (90)--(96);
\node[font=\small,anchor=west] at (7,0) {第 3 層}; \node[font=\small,anchor=west] at (7,-1) {第 2 層};
\node[font=\small,anchor=west] at (7,-2) {第 1 層};
\end{tikzpicture}
\end{center}
\begin{itemize}
\item (A) 錯。 根是 $79$。
\item (B) 對。 只有 $41 \to 61$ 一條水平連結（同層的父子），少於 $3$ 條。
\item (C) 錯。 $41$ 在第 $2$ 層，$49$ 在第 $1$ 層。 （換成看深度也不同：$41$ 深度 $1$，$49$ 深度 $3$。）
\item (D) 對。 $87$、$96$ 是 $90$ 的左右子節點，是兄弟。
\item (E) 錯。 $14$ 是 $41$ 的左子節點，$41$ 才是 $14$ 的父節點。
\end{itemize}
\end{pastsoln}
